% =============================================================================
% 01_prediction_design.tex — Problem 1: Prediction Design and Leakage Control
% =============================================================================
\section{Prediction Design and Leakage Control}
\label{sec:prediction-design}

% ---- 1.1 Target and Predictors ---------------------------------------------
\subsection{Target and Predictors}
\label{sec:target-predictors}

The \texttt{fat.csv} dataset contains $n = 252$ observations and 18 variables
describing body composition and anthropometric measurements of adult males.
The response variable is \texttt{brozek}, the body-fat percentage estimated via
Brozek's equation.

After inspection, the following variables are \textbf{excluded} from the
predictor set to prevent data leakage:
%
\begin{itemize}
    \item \texttt{siri} — an alternative body-fat estimate; including it would
          constitute a near-perfect proxy for the response. Both Siri and Brozek equations are derived directly from body density, making them fundamentally the same measure.
    \item \texttt{density} — body density is the direct input to both the Brozek
          and Siri formulas (e.g., $\text{Brozek} = 457 / \text{density} - 414.2$). Using it makes prediction trivially deterministic.
    \item \texttt{free} — fat-free weight is algebraically derived from body fat
          and total weight ($\text{free} = \text{weight} \times (1 - \text{body fat \%} / 100)$), creating a circular dependency and data leakage.
\end{itemize}

This leaves $p = 14$ candidate predictors (age, weight, height, adipos, neck,
chest, abdom, hip, thigh, knee, ankle, biceps, forearm, wrist).

\paragraph{Summary statistics.}
Table~\ref{tab:summary-stats} reports the descriptive statistics for all
retained predictors computed on the \emph{training set only}.

\begin{table}[H]
\centering
\caption{Descriptive statistics of retained predictors (training set).}
\label{tab:summary-stats}
\input{../output/tables/tab_p1_summary}
\end{table}

% ---- 1.2 Split and Train-Only Preprocessing --------------------------------
\subsection{Split and Train-Only Preprocessing}
\label{sec:split-preprocess}

The data are split into a training set ($n_{\mathrm{train}} \approx 201$,
$\sim 80\%$) and a held-out test set ($n_{\mathrm{test}} \approx 51$,
$\sim 20\%$) using a fixed random seed for reproducibility.

The exact fixed random seed used for this split was \texttt{240201}.

\paragraph{Standardisation.}
All predictors are centred and scaled using the \emph{training-set} mean and
standard deviation.  The same transformation is then applied to the test set to
prevent information leakage from the test data into the modelling pipeline.
Formally, for each predictor $j$:
\[
    \tilde{x}_{ij}
    = \frac{x_{ij} - \bar{x}_j^{\mathrm{train}}}
           {s_j^{\mathrm{train}}},
    \qquad i = 1, \dots, n.
\]

\paragraph{Cross-validation.}
We use $K = 10$-fold cross-validation (CV) throughout, with folds assigned once
and reused across all models to ensure fair comparison.

To strictly prevent data leakage, the test set labels are completely held out from model training and tuning. All standardisation is computed exclusively on the training set. We do not standardise within each CV fold since the primary pipeline standardises the entire training block prior to CV, which ensures the same preprocessing statistics are applied consistently across all folds.

% ---- 1.3 Training-Data Exploration -----------------------------------------
\subsection{Training-Data Exploration}
\label{sec:eda}

% -- Correlation heatmap
\includefigure{0.85\textwidth}{fig_corr_heatmap}%
    {Pearson correlation heatmap of training-set predictors and response.}%
    {fig:corr-heatmap}

The heatmap reveals that several anthropometric measurements are highly correlated with one another, such as \texttt{abdom}, \texttt{chest}, \texttt{hip}, and \texttt{weight}. High collinearity among these predictors can cause ordinary least squares estimates to become unstable. This motivates the use of regularisation techniques like Ridge and Lasso, which are designed to handle correlated features robustly.

% -- Response distribution
\includefigure{0.65\textwidth}{fig_p1_dist}%
    {Distribution of the response variable \texttt{brozek} in the training set.}%
    {fig:response-dist}

The distribution of the response variable \texttt{brozek} in the training set is roughly symmetric and bell-shaped, peaking around the mean of 18.67\%. It does not exhibit severe skewness, suggesting that no transformation of the response is strictly necessary for linear modelling.

% -- Boxplots of predictors
\includefigure{0.85\textwidth}{fig_p1_boxplots}%
    {Side-by-side boxplots of standardised predictors (training set).}%
    {fig:boxplots}

The side-by-side boxplots of the standardised predictors show varying spreads and several outliers, notably high-end outliers in the \texttt{weight}, \texttt{adipos}, and \texttt{ankle} variables. These extreme values could potentially have high leverage in a standard linear regression model.

% -- Pairwise scatter
\includefigure{0.95\textwidth}{fig_p1_pairwise}%
    {Pairwise scatter plots of selected predictors against \texttt{brozek}.}%
    {fig:pairwise}

The pairwise scatter plots show a generally linear and positive association between the response and the top four correlated predictors (\texttt{abdom}, \texttt{adipos}, \texttt{chest}, and \texttt{hip}). The relationships appear linear without obvious curvilinear patterns, supporting the use of linear models without polynomial transformations. The substantial multicollinearity between these top predictors strongly motivates the use of penalised models like Ridge and Lasso to regularise their coefficient estimates.
